% Frontmatter: Title, Abstract, Keywords
\title{\Huge\textbf{Comprehensive Catalog of Research Questions in Deep Learning for Skin Lesion Analysis}\\[0.3em]
\Large\textbf{A Systematic Methodological Extraction and Empirical Synthesis Across 39 Verified Peer-Reviewed Publications}}

\author{
    \textbf{Research Methodology Program (BCA / Computer Science)}\\[0.2em]
    \textit{Department of Computer Applications \& Medical Machine Learning Research Group}\\
    Course Code: BCA 7th-Semester --- Advanced Research Methodology\\
    Reference Repository: \texttt{downloaded\_journals/} (39 Verified Publications)
}

\date{\textbf{Academic Year: 2026} \quad | \quad Compilation Date: \today}

\maketitle
\thispagestyle{empty}

\begin{abstract}
\noindent Formulating clear, actionable, and theoretically grounded Research Questions (RQs) represents the foundational core of scientific inquiry in computer science, medical image analysis, and applied artificial intelligence. This monograph presents a comprehensive, systematic extraction and synthesis of the research questions investigated across \textbf{39 peer-reviewed research papers} (published between 2019 and 2026 in IEEE Access, IEEE Transactions, and premier medical informatics journals) centered on deep learning for skin cancer and dermatological lesion analysis. Following a rigorous evidence-bound methodology, we analyze both explicit research questions (notably from Systematic Literature Reviews adhering to the Kitchenham protocol) and core empirical questions formalized across four standardized academic dimensions: (1)~\textit{Architectural Innovation and Model Design}, (2)~\textit{Preprocessing, Feature Representation, and Artifact Mitigation}, (3)~\textit{Clinical Diagnostic Performance and Benchmarking}, and (4)~\textit{Explainability, Computational Efficiency, and Point-of-Care Translation}. For each paper, we document the precise methodological strategy employed, the quantitative empirical answers discovered, and the unresolved research gaps articulated by the authors. Furthermore, this catalog synthesizes overarching thematic paradigms across the 39 papers, including the architectural migration from early regularized CNNs toward Vision Transformers and Mixture-of-LoRA Experts, the resolution of extreme benchmark class imbalance, and the paradigm shift from post-hoc saliency heatmaps to neuro-symbolic clinical knowledge graphs and vision-language report generation. This work provides an exhaustive, publication-grade reference for researchers, graduate students, and clinical AI practitioners.
\end{abstract}

\begin{IEEEkeywords}
Research Methodology, Research Questions, Skin Cancer Detection, Deep Learning, Dermoscopy, Convolutional Neural Networks, Vision Transformers, Explainable AI, Systematic Literature Review, Evidence-Based Computing.
\end{IEEEkeywords}

\newpage
